\documentclass[a4paper,11pt]{ctexart}
\usepackage{amsmath, amssymb, amsfonts}
\usepackage{geometry}
\geometry{left=1.5cm, right=1.5cm, top=2.0cm, bottom=2.0cm, headsep=0.3cm, footskip=1cm}
\usepackage{listings}
\usepackage{xcolor}
\usepackage{tikz}
\usepackage{pgfplots}
\usepackage{hyperref}
\usepackage{fancyhdr}
\usepackage{tcolorbox}
\tcbuselibrary{breakable, skins}
\usepackage{array}
\usepackage{booktabs}
\usepackage[shortlabels]{enumitem}
\usepackage{needspace}
\usepackage{multirow}

\AtBeginDocument{
  \hypersetup{colorlinks=true, linkcolor=blue, filecolor=magenta, urlcolor=cyan}
}

\setlist{nosep, itemsep=0.8ex, parsep=0.1em}
\setlist[itemize]{leftmargin=1.8em}
\setlist[enumerate]{leftmargin=1.8em}

% ===== 颜色定义 =====
\definecolor{codegreen}{rgb}{0,0.6,0}
\definecolor{codegray}{rgb}{0.5,0.5,0.5}
\definecolor{codepurple}{rgb}{0.58,0,0.82}
\definecolor{codebg}{rgb}{0.98,0.98,0.98}
\definecolor{tipsblue}{rgb}{0.85,0.93,1.0}
\definecolor{highlightyellow}{rgb}{1.0,0.97,0.8}

% ===== 代码块样式 =====
\lstdefinestyle{mystyle}{
    backgroundcolor=\color{codebg},
    commentstyle=\color{codegreen},
    keywordstyle=\color{magenta}\bfseries,
    numberstyle=\tiny\color{codegray},
    stringstyle=\color{codepurple},
    basicstyle=\ttfamily\small,
    breakatwhitespace=false,
    breaklines=true,
    captionpos=b,
    keepspaces=true,
    numbers=left,
    numbersep=6pt,
    showspaces=false,
    showstringspaces=false,
    showtabs=false,
    tabsize=2,
    language=Python,
    frame=single,
    framerule=0.5pt,
    rulecolor=\color{gray!40}
}
\lstset{style=mystyle}

% ===== 彩色框 =====
\newtcolorbox{problembox}[1][]{
    colback=white,
    colframe=blue!60!black,
    title=\textbf{#1},
    fonttitle=\bfseries\small,
    boxrule=1pt,
    arc=3pt, outer arc=3pt,
    coltitle=black,
    before skip=10pt, after skip=8pt
}

\newtcolorbox{solutionbox}[1][]{
    enhanced, breakable,
    colback=green!3!white,
    colframe=green!50!black,
    title=\textbf{#1},
    fonttitle=\bfseries\small,
    boxrule=1pt,
    arc=3pt, outer arc=3pt,
    coltitle=black,
    before skip=8pt, after skip=10pt
}

\newtcolorbox{metaphorbox}[1][]{
    colback=orange!5!white,
    colframe=orange!70!black,
    title=\textbf{#1},
    fonttitle=\bfseries\small,
    boxrule=1pt,
    arc=3pt, outer arc=3pt,
    coltitle=black,
    before skip=8pt, after skip=8pt
}

\newtcolorbox{beginnerbox}[1][]{
    colback=tipsblue,
    colframe=blue!50!black,
    title=\textbf{#1},
    fonttitle=\bfseries\small,
    boxrule=1pt,
    arc=3pt, outer arc=3pt,
    coltitle=black,
    before skip=8pt, after skip=8pt
}

\newtcolorbox{appbox}[1][]{
    colback=green!3!white,
    colframe=green!60!black,
    title=\textbf{#1},
    fonttitle=\bfseries\small,
    boxrule=1pt,
    arc=3pt, outer arc=3pt,
    coltitle=black,
    before skip=8pt, after skip=10pt
}

\newtcolorbox{keybox}[1][]{
    colback=highlightyellow,
    colframe=orange!80!black,
    title=\textbf{#1},
    fonttitle=\bfseries\small,
    boxrule=1pt,
    arc=3pt, outer arc=3pt,
    coltitle=black,
    before skip=8pt, after skip=8pt
}

% ===== 页眉页脚 =====
\pagestyle{fancy}
\fancyhf{}
\fancyhead[R]{深度学习-第7章}
\fancyhead[L]{用 PyTorch 进行深度学习 · 练习题}
\fancyfoot[C]{\thepage}
\addtolength{\headheight}{2pt}

\parskip=2pt plus 1pt minus 0.5pt
\linespread{1.35}
\raggedbottom
\widowpenalty=10000
\clubpenalty=10000

\title{\textbf{深度学习\\第7章\quad 用 PyTorch 进行深度学习}\\[0.5em]
{\large\color{gray}——从激活函数到优化器，把神经网络的每个零件搞清楚}}
\date{\today}

\begin{document}

\maketitle

% ===== 章节导读 =====
\begin{beginnerbox}[本章题目导览：你将挑战哪些核心问题？]
    搭建神经网络就像组装一台精密机器：需要\textbf{激活函数}为输出引入非线性，需要\textbf{初始化策略}让训练一开始就走对方向，需要\textbf{损失函数}衡量预测有多准，还需要\textbf{优化器}一步步调整参数。本章 6 道题从零打通 PyTorch 深度学习的核心脉络：

    \begin{itemize}
        \item \textbf{Q7-01（辨析填空）}：四种激活函数横向对比——Sigmoid、Tanh、ReLU、Softmax，各自的公式、优缺点与适用场景是什么？
        \item \textbf{Q7-02（概念理解）}：参数初始化的重要性——为什么不能把权重全初始化为0？Xavier 和 Kaiming 初始化解决了什么问题？
        \item \textbf{Q7-03（计算题）}：损失函数手动推导——给定预测值和标签，分别计算交叉熵损失和 MSE 损失，把公式变成具体数字。
        \item \textbf{Q7-04（对比分析）}：五种优化器对比——Momentum、AdaGrad、RMSProp、Adam 各自的更新规则和核心思想是什么？
        \item \textbf{Q7-05（代码阅读）}：自定义模型代码追踪——\texttt{nn.Module} 子类的数据流形状逐层追踪，\texttt{forward} 方法怎么运作？
        \item \textbf{Q7-06（综合应用）}：房价预测模型全流程——特征工程、模型搭建、损失函数选择和训练过程的完整实现思路。
    \end{itemize}
\end{beginnerbox}

\section{第 7 章\quad 用 PyTorch 进行深度学习}

在前几章，我们已经理解了神经网络的基本原理：数据流入，经过若干层变换，输出预测结果，和标签对比，反向传播更新权重。

但"知道原理"和"真正能搭"之间，还差很多零件：

\begin{itemize}
    \item 每一层之后要不要加激活函数？加哪种？
    \item 权重一开始怎么设置？全设为 0 行不行？
    \item 预测值和真实值的"差距"用什么公式来量化？
    \item 梯度算完了，用什么策略更新参数最有效？
\end{itemize}

这些问题，就是本章要解决的。\textbf{每一个零件都不是装饰品——选错了，模型可能根本训练不起来。}

\begin{metaphorbox}[先建立一个总感觉：PyTorch 搭网络像什么？]
    \textbf{把搭建神经网络比作开一家餐厅。}

    \begin{itemize}
        \item \textbf{激活函数} = 厨师的调味料：没有调味料，所有菜都是白水煮——网络没有激活函数，所有层叠加起来还是一个线性变换，表达能力极差；
        \item \textbf{参数初始化} = 餐厅开业前的备货：备太少（权重趋于0）厨师无从下手，备太多（权重过大）食材腐坏——初始化决定训练能不能顺利起跑；
        \item \textbf{损失函数} = 顾客的评分系统：没有评分，厨师不知道菜好不好；评分维度选错，厨师改进方向也就跑偏了；
        \item \textbf{优化器} = 根据评分改进菜谱的策略：是只看今天的差评（SGD），还是记住历史教训（Momentum），还是对每道菜用不同力度改进（Adam）？
    \end{itemize}
\end{metaphorbox}


%% ============================================================
\Needspace{5\baselineskip}
\begin{problembox}[Q7-01\quad 四种激活函数横向对比]
    \textbf{题目描述：}\\
    文档第 7.1 节介绍了神经网络中常用的四种激活函数：Sigmoid、Tanh、ReLU 和 Softmax。在做题之前，先来认识一下它们的本质区别。

    \vspace{0.3cm}
    \textbf{问题：}
    \begin{enumerate}[(1)]
        \item 填写下表，将四种激活函数的关键特征补全：

        \begin{center}
        \begin{tabular}{p{2.0cm} p{4.0cm} p{3.5cm} p{3.5cm}}
            \toprule
            激活函数 & 数学公式 & 主要缺点 & 典型应用场景 \\
            \midrule
            Sigmoid & $\sigma(x) = \dfrac{1}{1+e^{-x}}$ & 存在\underline{\hspace{1.5cm}}问题，输出非零均值 & \underline{\hspace{2.5cm}} \\[2.0em]
            Tanh & $\tanh(x) = \dfrac{e^x - e^{-x}}{e^x + e^{-x}}$ & 仍存在\underline{\hspace{1.0cm}}问题 & 输出需要零均值的\underline{\hspace{1.5cm}} \\[2.0em]
            ReLU & $\text{ReLU}(x) = \underline{\hspace{2.5cm}}$ & 负输入区域存在\underline{\hspace{1.5cm}}现象 & 深层网络的\underline{\hspace{1.5cm}}层 \\[2.0em]
            Softmax & $p_i = \dfrac{e^{z_i}}{\sum_j e^{z_j}}$ & 输出之和为1，适合\underline{\hspace{1.5cm}} & 多分类任务的\underline{\hspace{1.5cm}}层 \\
            \bottomrule
        \end{tabular}
        \end{center}

        \item "梯度消失"（Vanishing Gradient）是什么？为什么 Sigmoid 和 Tanh 在深层网络中容易引发这个问题？
        \item ReLU 的全称是什么？它相比 Sigmoid/Tanh 的最大优势是什么（从计算效率和梯度角度回答）？
    \end{enumerate}

    \vspace{0.2cm}
    {\color{gray}\small\textit{来源溯源：[文档第 7 章 7.1 激活函数]}}
\end{problembox}

\begin{beginnerbox}[小白必读：为什么神经网络必须要激活函数？]
    \begin{itemize}
        \item \textbf{没有激活函数会发生什么？}假设有三层线性变换：$y = W_3(W_2(W_1 x))$。由于矩阵乘法满足结合律，这等价于 $y = (W_3 W_2 W_1) x$，本质上还是一个线性变换！无论网络多深，表达能力和单层一模一样，完全无法拟合非线性的现实数据。
        \item \textbf{激活函数的本质}：在每一层输出后施加一个非线性变换，打破线性叠加的束缚，让网络有能力学习任意复杂的函数。
        \item \textbf{不同位置用不同激活函数}：隐藏层通常用 ReLU（计算快，梯度稳定）；二分类输出层用 Sigmoid（输出0到1之间的概率）；多分类输出层用 Softmax（输出加起来等于1的概率分布）；回归输出层通常不加激活函数（输出任意实数）。
    \end{itemize}
\end{beginnerbox}

\begin{metaphorbox}[生活比喻：激活函数像什么？]
    想象你在做一道判断题测验：

    \begin{itemize}
        \item \textbf{Sigmoid}：把任何分数都压缩到 $(0, 1)$ 之间——"你有多大把握这道题是对的？"输出越接近 1，越有把握。但当分数特别高或特别低时，Sigmoid 的输出变化极小，梯度几乎为0，反向传播时信号就"消失"了，深层网络因此难以更新。
        \item \textbf{ReLU}：正分保留，零分或负分直接归零——"这个特征有用就留，没用就扔"。简单粗暴，但计算极快，而且正值区域梯度恒为1，信号可以畅通地传到最底层。
        \item \textbf{Softmax}：把一组原始分数变成"投票比例"——"三个候选答案，分别有82\%、11\%、7\% 的概率是对的"。结果之和必须等于1，天然适合多分类。
    \end{itemize}
\end{metaphorbox}

\begin{solutionbox}[参考答案与详细解析]
    \textbf{(1) 激活函数对比表（填写内容）：}

    \begin{center}
    \begin{tabular}{p{2.0cm} p{4.0cm} p{3.5cm} p{3.5cm}}
        \toprule
        激活函数 & 数学公式 & 主要缺点 & 典型应用场景 \\
        \midrule
        Sigmoid & $\sigma(x)=\dfrac{1}{1+e^{-x}}$ & 梯度消失，输出非零均值 & 二分类输出层 \\[1.5em]
        Tanh & $\tanh(x)=\dfrac{e^x-e^{-x}}{e^x+e^{-x}}$ & 仍存在梯度消失问题 & 输出需要零均值的隐藏层 \\[1.5em]
        ReLU & $\max(0, x)$ & 负输入区域存在神经元死亡现象 & 深层网络的隐藏层 \\[1.5em]
        Softmax & $p_i=\dfrac{e^{z_i}}{\sum_j e^{z_j}}$ & 输出之和为1，适合互斥分类 & 多分类任务的输出层 \\
        \bottomrule
    \end{tabular}
    \end{center}

    \textbf{(2) 梯度消失原因：}

    Sigmoid 函数的导数为 $\sigma'(x) = \sigma(x)(1-\sigma(x))$，最大值仅为 $0.25$（在 $x=0$ 时取得）。Tanh 的导数最大值为 $1$（在 $x=0$ 时），远离原点后也迅速趋近于0。

    在反向传播中，梯度是逐层相乘的。若每层梯度最大为 $0.25$，经过 10 层后，梯度缩小至 $0.25^{10} \approx 10^{-6}$，几乎为0，底层网络参数无法有效更新。

    \textbf{(3) ReLU 的优势：}

    ReLU（Rectified Linear Unit，修正线性单元）的公式为 $\max(0,x)$，优势体现在两方面：
    \begin{itemize}
        \item \textbf{计算极快}：只需判断正负，无需计算指数，比 Sigmoid/Tanh 快数倍；
        \item \textbf{梯度稳定}：正值区域梯度恒为1，不会随层数加深而衰减，有效缓解梯度消失。
    \end{itemize}
\end{solutionbox}


%% ============================================================
\Needspace{5\baselineskip}
\begin{problembox}[Q7-02\quad 参数初始化：为什么不能全设为 0？]
    \textbf{题目描述：}\\
    文档第 7.2 节介绍了多种参数初始化方法。一个常见的误区是"把所有权重初始化为0，反正训练会自动调整"。

    \vspace{0.3cm}
    \textbf{问题：}
    \begin{enumerate}[(1)]
        \item \textbf{对称性破坏问题}：如果把一个全连接层（\texttt{nn.Linear}）的所有权重初始化为相同的值（包括全0），会出现什么问题？从反向传播的角度解释"对称性"为什么无法被打破。

        \item \textbf{Xavier 初始化}：适用于 Sigmoid/Tanh 激活函数之前的层，权重从以下范围均匀采样：
        \[
        W \sim \mathcal{U}\!\left(-\frac{\sqrt{6}}{\sqrt{n_{\text{in}}+n_{\text{out}}}},\ \frac{\sqrt{6}}{\sqrt{n_{\text{in}}+n_{\text{out}}}}\right)
        \]
        若一个线性层输入特征数 $n_{\text{in}}=128$，输出特征数 $n_{\text{out}}=64$，计算 Xavier 初始化的采样范围上界（保留两位小数）。

        \item \textbf{Kaiming 初始化}：专为 ReLU 激活函数设计，权重从均值为0、标准差为 $\sqrt{2/n_{\text{in}}}$ 的正态分布中采样。解释为什么 ReLU 需要专门的初始化策略（提示：ReLU 会"杀死"一半神经元）。

        \item 在 PyTorch 中，\texttt{nn.Linear} 层默认使用哪种初始化方式？如何手动将某层权重改为 Kaiming 正态初始化？写出相关代码。
    \end{enumerate}

    \vspace{0.2cm}
    {\color{gray}\small\textit{来源溯源：[文档第 7 章 7.2 参数初始化和正则化]}}
\end{problembox}

\begin{beginnerbox}[小白必读：初始化为什么这么重要？]
    \begin{itemize}
        \item \textbf{权重不是随便设的}：如果初始化过大，激活值和梯度会爆炸（梯度爆炸）；如果初始化过小，激活值趋近于0，梯度也接近0（梯度消失）。两种情况都会让训练无法有效进行。
        \item \textbf{全0初始化的致命问题}：如果同一层所有神经元的权重完全相同，那么正向传播时它们的输出完全相同，反向传播时它们接收到的梯度也完全相同，更新后仍然完全相同——这层神经元永远无法学到不同的特征，相当于只有一个神经元在工作。
        \item \textbf{为什么需要随机初始化}：随机初始化打破了对称性，让每个神经元从一开始就走向不同的"方向"，学到不同的特征。
    \end{itemize}
\end{beginnerbox}

\begin{solutionbox}[参考答案与详细解析]
    \textbf{(1) 对称性问题：}

    若所有权重 $w_{ij}^{(l)} = c$（常数），则同一层的每个神经元接收到完全相同的输入，产生完全相同的输出（正向传播对称）。反向传播时，损失对该层每个神经元的梯度也完全相同，参数更新幅度一致，导致更新后权重仍然相同——\textbf{对称性永远无法被打破}，整个层退化为"一个神经元"，无论宽度多大都形同虚设。

    \vspace{0.5em}
    \textbf{(2) Xavier 初始化上界计算：}

    \[
    \text{上界} = \frac{\sqrt{6}}{\sqrt{n_{\text{in}}+n_{\text{out}}}} = \frac{\sqrt{6}}{\sqrt{128+64}} = \frac{\sqrt{6}}{\sqrt{192}} = \frac{2.449}{13.856} \approx \boxed{0.18}
    \]

    采样范围为 $[-0.18,\ 0.18]$。

    \vspace{0.5em}
    \textbf{(3) ReLU 需要专门初始化的原因：}

    ReLU 对所有负输入输出0，相当于每次前向传播中约一半的神经元被"关闭"（输出为0，梯度也为0）。若使用 Xavier 初始化（针对输出方差设计），实际上只有约一半的神经元在传递信号，等效扇入减半，方差会缩小约2倍。Kaiming 初始化将标准差设为 $\sqrt{2/n_{\text{in}}}$（比 Xavier 多了因子 $\sqrt{2}$），恰好补偿了 ReLU 使方差减半的效果，使各层激活值的方差保持稳定。

    \vspace{0.5em}
    \textbf{(4) PyTorch 代码：}

    \texttt{nn.Linear} 默认使用 Kaiming 均匀初始化（\texttt{kaiming\_uniform\_}）。手动改为 Kaiming 正态初始化：

    \begin{lstlisting}[language=Python]
import torch.nn as nn

layer = nn.Linear(128, 64)

# 手动将权重改为 Kaiming 正态初始化（适合 ReLU）
nn.init.kaiming_normal_(layer.weight, mode='fan_in',
                         nonlinearity='relu')
# 偏置通常初始化为 0
nn.init.zeros_(layer.bias)
    \end{lstlisting}

    \begin{keybox}[常用初始化方法速查]
        \begin{itemize}
            \item \textbf{全零初始化}：\texttt{nn.init.zeros\_(tensor)} —— 仅用于偏置，权重千万不要用
            \item \textbf{Xavier 均匀}：\texttt{nn.init.xavier\_uniform\_(tensor)} —— 配合 Sigmoid/Tanh
            \item \textbf{Kaiming 正态}：\texttt{nn.init.kaiming\_normal\_(tensor, nonlinearity='relu')} —— 配合 ReLU
            \item \textbf{正态分布}：\texttt{nn.init.normal\_(tensor, mean=0, std=0.01)} —— 通用小方差初始化
        \end{itemize}
    \end{keybox}
\end{solutionbox}


%% ============================================================
\Needspace{5\baselineskip}
\begin{problembox}[Q7-03\quad 损失函数手动计算：交叉熵 vs MSE]
    \textbf{题目描述：}\\
    文档第 7.4 节介绍了分类任务和回归任务的常用损失函数。损失函数是训练信号的来源——选错了，模型会朝错误的方向优化。

    \vspace{0.3cm}
    \textbf{场景一：多分类任务（交叉熵损失）}

    一个三分类模型，对某样本输出的 logits 为 $[2.0,\ 1.0,\ 0.1]$，真实标签为第 0 类（即 $y=0$）。

    \vspace{0.2cm}
    \textbf{场景二：回归任务（MSE 损失）}

    一个房价预测模型，对 4 个样本的预测值（万元）为 $[\hat{y}_1, \hat{y}_2, \hat{y}_3, \hat{y}_4] = [102, 95, 310, 88]$，真实值为 $[100, 98, 300, 90]$。

    \vspace{0.3cm}
    \textbf{问题：}
    \begin{enumerate}[(1)]
        \item 对场景一，先计算 softmax 概率，再计算交叉熵损失：
        \[
        \mathcal{L}_{\text{CE}} = -\log(p_{y})
        \]
        其中 $p_y$ 是真实类别对应的预测概率。（结果保留两位小数）

        \item 对场景二，计算均方误差（MSE）损失：
        \[
        \mathcal{L}_{\text{MSE}} = \frac{1}{N}\sum_{i=1}^{N}(\hat{y}_i - y_i)^2
        \]
        （结果保留两位小数）

        \item PyTorch 的 \texttt{nn.CrossEntropyLoss} 内部已经包含了 Softmax 操作。如果在模型的输出层\textbf{已经手动加了 Softmax}，再使用 \texttt{nn.CrossEntropyLoss}，会产生什么问题？应该怎么修改？

        \item 为什么回归任务不能用交叉熵损失，分类任务不应该用 MSE 损失？各举一个具体理由。
    \end{enumerate}

    \vspace{0.2cm}
    {\color{gray}\small\textit{来源溯源：[文档第 7 章 7.4 损失函数]}}
\end{problembox}

\begin{beginnerbox}[小白必读：损失函数是什么？为什么要分任务类型？]
    \begin{itemize}
        \item \textbf{损失函数的本质}：量化"预测值和真实值有多不一样"的函数。训练就是通过梯度下降不断最小化损失函数的过程。
        \item \textbf{为什么要区分分类和回归？}分类任务的标签是离散类别（猫/狗/鸟），预测的是概率分布；回归任务的标签是连续数值（房价/温度），预测的是具体数值。两种情况下"误差"的计量方式完全不同。
        \item \textbf{CrossEntropyLoss 一定要配 logits 用}：PyTorch 的 \texttt{nn.CrossEntropyLoss} 内置了 softmax + log + 取均值的完整计算流程，输入必须是未归一化的原始分数（logits），而不是经过 softmax 后的概率。如果先 softmax 再送进去，相当于做了两次 softmax，结果会严重错误。
    \end{itemize}
\end{beginnerbox}

\begin{solutionbox}[参考答案与详细解析]
    \textbf{(1) 交叉熵损失计算：}

    第一步，计算 softmax 概率（logits $= [2.0, 1.0, 0.1]$）：
    \[
    e^{2.0} \approx 7.389,\quad e^{1.0} \approx 2.718,\quad e^{0.1} \approx 1.105
    \]
    \[
    \text{总和} = 7.389 + 2.718 + 1.105 = 11.212
    \]
    \[
    p_0 = \frac{7.389}{11.212} \approx 0.659,\quad p_1 \approx 0.242,\quad p_2 \approx 0.099
    \]

    第二步，计算交叉熵损失（真实类别为第 0 类，$p_y = p_0 \approx 0.659$）：
    \[
    \mathcal{L}_{\text{CE}} = -\log(0.659) \approx \boxed{0.42}
    \]

    \textbf{(2) MSE 损失计算：}
    \[
    \mathcal{L}_{\text{MSE}} = \frac{(102-100)^2 + (95-98)^2 + (310-300)^2 + (88-90)^2}{4}
    = \frac{4 + 9 + 100 + 4}{4} = \frac{117}{4} = \boxed{29.25}
    \]

    \textbf{(3) 双重 Softmax 的问题：}

    \texttt{nn.CrossEntropyLoss} 内部等价于：$\mathcal{L} = -\log\!\left(\text{softmax}(\text{logits})_y\right)$。如果模型输出层已手动加了 softmax，那么传入的 $p = \text{softmax}(\text{logits})$，各元素在 $(0,1)$ 之间，\texttt{CrossEntropyLoss} 再对它做一次 softmax，会把概率值当 logits 处理，使得概率分布被极度压平，损失计算完全失真，训练无法正常进行。

    \textbf{修改方案}：
    \begin{itemize}
        \item \textbf{方案一（推荐）}：去掉模型输出层的 softmax，保持输出为 logits，直接配合 \texttt{nn.CrossEntropyLoss}；
        \item \textbf{方案二}：保留 softmax，将损失函数改为 \texttt{nn.NLLLoss}，并在输出后加 \texttt{torch.log}（因为 NLLLoss 期望输入 log-probabilities）。
    \end{itemize}

    \textbf{(4) 损失函数不可乱用的原因：}

    \begin{itemize}
        \item \textbf{回归不能用交叉熵}：交叉熵需要概率分布作为输入，而回归输出是任意实数，无法直接套用概率公式。强行使用会导致数学上无意义（对负数取对数会出错）。
        \item \textbf{分类不应用 MSE}：分类标签（如 $y \in \{0,1,2\}$）是无序的类别编号，MSE 会把类别间距离的大小当作重要信息处理，而实际上类别 0 和类别 2 并不比类别 0 和类别 1 "差距更大"。此外，MSE 在分类问题中梯度较小，训练收敛极慢。
    \end{itemize}
\end{solutionbox}


%% ============================================================
\Needspace{5\baselineskip}
\begin{problembox}[Q7-04\quad 五种优化器的核心思想对比]
    \textbf{题目描述：}\\
    文档第 7.5 节介绍了多种参数更新方法。优化器决定了"梯度算出来之后，怎么用它来更新参数"。选错了优化器，训练可能收敛极慢，或者在最优解附近震荡。

    \vspace{0.3cm}
    \textbf{问题：}
    \begin{enumerate}[(1)]
        \item 填写下表，将五种优化器的关键特征补全：

        \begin{center}
        \begin{tabular}{p{2.0cm} p{5.0cm} p{4.5cm}}
            \toprule
            优化器 & 核心更新规则（用文字描述） & 主要优缺点 \\
            \midrule
            SGD & 沿梯度负方向直接更新，步长为\underline{\hspace{1.0cm}} & 简单，但容易陷入\underline{\hspace{1.5cm}}，收敛慢 \\[1.5em]
            Momentum & 在 SGD 基础上累积\underline{\hspace{1.5cm}}，赋予参数"惯性" & 加速收敛，但\underline{\hspace{1.5cm}}参数较多 \\[1.5em]
            AdaGrad & 对每个参数使用\underline{\hspace{1.5cm}}学习率，历史梯度平方和大的参数学习率变\underline{\hspace{0.8cm}} & 自适应学习率，但学习率单调\underline{\hspace{1.0cm}}，后期更新极慢 \\[1.5em]
            RMSProp & 改进 AdaGrad，用\underline{\hspace{1.5cm}}平均代替全量累积 & 解决了学习率\underline{\hspace{1.5cm}}问题 \\[1.5em]
            Adam & 结合 Momentum 和 RMSProp，同时维护\underline{\hspace{1.5cm}}和\underline{\hspace{1.5cm}}的指数移动平均 & 效果好，但\underline{\hspace{1.5cm}}参数有时需调整 \\
            \bottomrule
        \end{tabular}
        \end{center}

        \item \textbf{学习率衰减}：训练过程中为什么要让学习率随时间减小？常见的衰减策略有哪些（至少列举两种）？

        \item \textbf{Adam 的偏差修正}：Adam 的动量估计在初始阶段会偏向于0，因此引入了偏差修正：
        \[
        \hat{m}_t = \frac{m_t}{1 - \beta_1^t}, \quad \hat{v}_t = \frac{v_t}{1 - \beta_2^t}
        \]
        解释为什么初始阶段 $m_t$ 会偏向于0（提示：$\beta_1$ 通常为 0.9，$m_0=0$，$m_1 = (1-\beta_1) g_1$）。

        \item 在 PyTorch 中，如何为网络不同层设置不同的学习率（例如：特征提取层学习率为 $10^{-4}$，分类头学习率为 $10^{-3}$）？写出关键代码片段。
    \end{enumerate}

    \vspace{0.2cm}
    {\color{gray}\small\textit{来源溯源：[文档第 7 章 7.5 参数更新方法]}}
\end{problembox}

\begin{beginnerbox}[小白必读：优化器是什么？为什么不直接用梯度？]
    \begin{itemize}
        \item \textbf{最简单的更新规则}：$w \leftarrow w - \eta \cdot g$，其中 $\eta$ 是学习率，$g$ 是梯度。这就是最基础的 SGD（随机梯度下降）。
        \item \textbf{为什么需要更复杂的优化器？}原始 SGD 有几个痛点：(1) 在平坦区域梯度很小，更新极慢；(2) 在陡峭区域梯度很大，更新容易震荡；(3) 所有参数共用同一个学习率，不够灵活。各种改进版优化器，就是针对这些痛点的解决方案。
        \item \textbf{Adam 是目前最常用的选择}：结合了动量（记住历史方向）和自适应学习率（对不同参数用不同步长），在绝大多数任务上开箱即用效果不错，是入门深度学习的默认选择。
    \end{itemize}
\end{beginnerbox}

\begin{metaphorbox}[生活比喻：五种优化器像五种下山策略]
    你在一座山上，要下到最低点（最优解），眼睛只能看到脚下的坡度（梯度）：

    \begin{itemize}
        \item \textbf{SGD}：每步严格沿当前最陡方向走固定步长，路线可能锯齿形剧烈震荡；
        \item \textbf{Momentum}：像滚下山的石头，有惯性，能越过小坑和轻微的鞍部，但可能冲过头；
        \item \textbf{AdaGrad}：在平坦路段迈大步，在陡峭路段迈小步——但走着走着，步子会越来越小，最终走不动；
        \item \textbf{RMSProp}：和 AdaGrad 类似，但"遗忘"很久以前的坡度，始终聚焦于近期信息，步子不会越来越小；
        \item \textbf{Adam}：既有滚石的惯性，又能根据近期坡度自动调节步长，是综合效果最好的"登山向导"。
    \end{itemize}
\end{metaphorbox}

\begin{solutionbox}[参考答案与详细解析]
    \textbf{(1) 优化器对比表（填写内容）：}

    \begin{center}
    \begin{tabular}{p{2.0cm} p{5.0cm} p{4.5cm}}
        \toprule
        优化器 & 核心更新规则 & 主要优缺点 \\
        \midrule
        SGD & 沿梯度负方向直接更新，步长为学习率 $\eta$ & 简单，但容易陷入局部极小值，收敛慢 \\[1em]
        Momentum & 累积历史梯度，赋予参数"惯性" & 加速收敛，但超参数较多 \\[1em]
        AdaGrad & 历史梯度平方和大的参数学习率变小（自适应） & 自适应，但学习率单调递减，后期更新极慢 \\[1em]
        RMSProp & 用指数移动平均代替全量累积 & 解决了学习率一直衰减问题 \\[1em]
        Adam & 同时维护一阶动量和二阶动量的指数移动平均 & 效果好，但默认超参数有时需调整 \\
        \bottomrule
    \end{tabular}
    \end{center}

    \textbf{(2) 学习率衰减策略：}

    训练初期需要大学习率以快速下降；训练后期接近最优解时，大学习率会导致在最优解附近震荡而无法收敛，需要小学习率精细调整。常见策略：
    \begin{itemize}
        \item \textbf{步进衰减（StepLR）}：每隔固定 epoch 数乘以衰减因子，如每 10 个 epoch 乘以 0.1；
        \item \textbf{余弦退火（CosineAnnealingLR）}：学习率按余弦函数平滑衰减，训练结束时接近0；
        \item \textbf{指数衰减（ExponentialLR）}：每个 epoch 乘以固定衰减系数 $\gamma$，如 $\eta_t = \eta_0 \cdot \gamma^t$。
    \end{itemize}

    \textbf{(3) Adam 偏差修正的原因：}

    $m_0 = 0$，$m_1 = \beta_1 m_0 + (1-\beta_1) g_1 = 0.1 \cdot g_1$。仅为真实梯度的 $10\%$，严重低估了真实动量大小。偏差修正 $\hat{m}_1 = m_1 / (1-\beta_1^1) = m_1 / 0.1 = g_1$，将其还原为真实梯度值。随着训练步数 $t$ 增大，$\beta_1^t \to 0$，修正项趋近于1，影响逐渐消失。

    \textbf{(4) 不同层设置不同学习率：}

    \begin{lstlisting}[language=Python]
import torch.optim as optim

model = MyModel()

# 分组设置不同学习率
optimizer = optim.Adam([
    {'params': model.feature_extractor.parameters(),
     'lr': 1e-4},
    {'params': model.classifier.parameters(),
     'lr': 1e-3}
], lr=1e-3)  # 默认学习率（作为 fallback）

# 之后正常训练即可
optimizer.zero_grad()
loss.backward()
optimizer.step()
    \end{lstlisting}
\end{solutionbox}


%% ============================================================
\Needspace{5\baselineskip}
\begin{problembox}[Q7-05\quad 自定义模型代码阅读：数据流形状追踪]
    \textbf{题目描述：}\\
    文档第 7.3 节介绍了如何用 \texttt{nn.Module} 搭建自定义神经网络。以下是一个用于图像二分类的模型代码：

    \begin{lstlisting}[language=Python]
import torch
import torch.nn as nn

class BinaryClassifier(nn.Module):
    def __init__(self, input_dim):
        super(BinaryClassifier, self).__init__()
        self.fc1 = nn.Linear(input_dim, 128)
        self.bn1 = nn.BatchNorm1d(128)
        self.dropout = nn.Dropout(p=0.3)
        self.fc2 = nn.Linear(128, 64)
        self.fc3 = nn.Linear(64, 1)
        self.relu = nn.ReLU()
        self.sigmoid = nn.Sigmoid()

    def forward(self, x):
        x = self.relu(self.bn1(self.fc1(x)))  # 第1步
        x = self.dropout(x)                   # 第2步
        x = self.relu(self.fc2(x))            # 第3步
        x = self.sigmoid(self.fc3(x))         # 第4步
        return x

model = BinaryClassifier(input_dim=784)
    \end{lstlisting}

    输入 \texttt{x} 的形状为 \texttt{[batch\_size=32, input\_dim=784]}（32 张展平的 28$\times$28 图像）。

    \vspace{0.2cm}
    \textbf{问题：}
    \begin{enumerate}[(1)]
        \item 逐步追踪各阶段张量的形状：
        \begin{enumerate}[a.]
            \item 模型输入 \texttt{x}：\underline{\hspace{4cm}}
            \item 第1步之后（\texttt{relu(bn1(fc1(x)))}）：\underline{\hspace{4cm}}
            \item 第2步之后（\texttt{dropout(x)}）：\underline{\hspace{4cm}}
            \item 第3步之后（\texttt{relu(fc2(x))}）：\underline{\hspace{4cm}}
            \item 第4步之后（模型输出）：\underline{\hspace{4cm}}
        \end{enumerate}

        \item \texttt{nn.BatchNorm1d(128)} 的作用是什么？为什么要把它放在线性层\textbf{之后}、激活函数\textbf{之前}？

        \item \texttt{nn.Dropout(p=0.3)} 在\textbf{训练阶段}和\textbf{推理阶段}的行为有什么不同？在 PyTorch 中如何切换这两种模式？

        \item 如果将这个模型改为用 \texttt{nn.Sequential} 的写法，请写出等价的代码（忽略 Dropout）。
    \end{enumerate}

    \vspace{0.2cm}
    {\color{gray}\small\textit{来源溯源：[文档第 7 章 7.3 搭建神经网络 / 7.2.8 Dropout]}}
\end{problembox}

\begin{beginnerbox}[小白必读：\texttt{nn.Module} 的两个核心方法]
    \begin{itemize}
        \item \textbf{\texttt{\_\_init\_\_}：定义零件}。在这里声明所有要用到的层（\texttt{nn.Linear}、\texttt{nn.ReLU} 等）。PyTorch 会自动追踪这些层的参数，让优化器能找到它们。
        \item \textbf{\texttt{forward}：定义数据怎么流动}。调用 \texttt{model(x)} 时，实际上执行的就是 \texttt{forward(x)}。层的定义顺序（\texttt{\_\_init\_\_}）和使用顺序（\texttt{forward}）不必相同，这是 \texttt{nn.Module} 比 \texttt{nn.Sequential} 灵活的地方——可以实现跳跃连接、多输入、多输出等复杂结构。
        \item \textbf{不要忘记 \texttt{super().\_\_init\_\_()}}：继承 \texttt{nn.Module} 时，第一行必须调用父类初始化，否则参数注册机制会失效，优化器找不到参数。
    \end{itemize}
\end{beginnerbox}

\begin{solutionbox}[参考答案与详细解析]
    \textbf{(1) 逐层形状追踪：}

    \begin{center}
    \begin{tabular}{clll}
        \toprule
        步骤 & 操作 & 输出形状 & 说明 \\
        \midrule
        输入 & \texttt{x} & \texttt{[32, 784]} & 32张展平的图像 \\
        第1步 & \texttt{relu(bn1(fc1(x)))} & \texttt{[32, 128]} & Linear(784→128)，BN，ReLU \\
        第2步 & \texttt{dropout(x)} & \texttt{[32, 128]} & Dropout 不改变形状 \\
        第3步 & \texttt{relu(fc2(x))} & \texttt{[32, 64]} & Linear(128→64)，ReLU \\
        第4步 & \texttt{sigmoid(fc3(x))} & \texttt{[32, 1]} & Linear(64→1)，Sigmoid \\
        \bottomrule
    \end{tabular}
    \end{center}

    \textbf{(2) BatchNorm1d 的作用及位置：}

    批归一化（Batch Normalization）对每一批数据的每个特征维度独立进行归一化（减去均值，除以标准差），使激活值保持在稳定的范围内，有助于加速训练、缓解梯度问题，并起到一定的正则化作用。

    放在线性层\textbf{之后}、激活函数\textbf{之前}，是因为：归一化的对象是线性变换的输出（未激活的值），使其接近标准正态分布，再经过激活函数时梯度信号最稳定。若放在激活函数之后，ReLU 已经截断了负值，归一化效果会大打折扣。

    \textbf{(3) Dropout 在训练 vs 推理阶段的差异：}

    \begin{itemize}
        \item \textbf{训练阶段}：以概率 $p=0.3$ 随机将 30\% 的神经元输出置为0，迫使网络不依赖某几个特定神经元，从而起到正则化作用。未被置零的神经元输出会乘以 $\frac{1}{1-p}$ 进行缩放，保持期望值不变。
        \item \textbf{推理阶段}：所有神经元正常输出，不做任何置零操作（因为推理时需要稳定的确定性输出）。
    \end{itemize}

    在 PyTorch 中通过以下方式切换：
    \begin{lstlisting}[language=Python]
model.train()   # 切换为训练模式（Dropout 生效）
model.eval()    # 切换为推理模式（Dropout 关闭）
    \end{lstlisting}

    \textbf{(4) 等价的 Sequential 写法：}

    \begin{lstlisting}[language=Python]
model = nn.Sequential(
    nn.Linear(784, 128),
    nn.BatchNorm1d(128),
    nn.ReLU(),
    nn.Linear(128, 64),
    nn.ReLU(),
    nn.Linear(64, 1),
    nn.Sigmoid()
)
    \end{lstlisting}

    \begin{keybox}[\texttt{nn.Module} vs \texttt{nn.Sequential} 的选择建议]
        \begin{itemize}
            \item \textbf{用 Sequential}：结构简单，数据单路径从头流到尾，代码更简洁；
            \item \textbf{用 Module 子类}：需要跳跃连接（ResNet）、多分支（Inception）、多输入或多输出，或者 forward 中有条件逻辑时，必须用 Module 子类。
        \end{itemize}
    \end{keybox}
\end{solutionbox}


%% ============================================================
\Needspace{5\baselineskip}
\begin{problembox}[Q7-06\quad 综合应用：房价预测全流程分析]
    \textbf{题目描述：}\\
    文档第 7.6 节给出了一个用神经网络预测房价的完整案例。以下是简化后的核心代码骨架：

    \begin{lstlisting}[language=Python]
import torch
import torch.nn as nn
import pandas as pd
from sklearn.preprocessing import StandardScaler

# ---- 7.6.2 特征工程 ----
df = pd.read_csv('house_prices.csv')
# 处理缺失值：数值列用均值填充，类别列用众数填充
num_cols = df.select_dtypes(include='number').columns
cat_cols = df.select_dtypes(exclude='number').columns
df[num_cols] = df[num_cols].fillna(df[num_cols].mean())
df[cat_cols] = df[cat_cols].fillna(df[cat_cols].mode().iloc[0])
# 类别列独热编码
df = pd.get_dummies(df, columns=cat_cols)
# 特征标准化
scaler = StandardScaler()
X = scaler.fit_transform(df.drop('SalePrice', axis=1))
y = df['SalePrice'].values

# ---- 7.6.3 搭建模型 ----
class HousePriceModel(nn.Module):
    def __init__(self, input_dim):
        super().__init__()
        self.net = nn.Sequential(
            nn.Linear(input_dim, 256),
            nn.ReLU(),
            nn.Dropout(0.2),
            nn.Linear(256, 128),
            nn.ReLU(),
            nn.Linear(128, 1)         # 输出层无激活函数
        )

    def forward(self, x):
        return self.net(x)

# ---- 7.6.4 损失函数 ----
loss_fn = nn.MSELoss()

# ---- 7.6.5 模型训练 ----
optimizer = torch.optim.Adam(model.parameters(), lr=1e-3)
    \end{lstlisting}

    \vspace{0.2cm}
    \textbf{问题：}
    \begin{enumerate}[(1)]
        \item \textbf{特征工程}：代码中对类别特征进行了"独热编码"（One-Hot Encoding）。解释独热编码的原理，并说明为什么不能直接用类别的整数编码（如将"北京=1，上海=2，广州=3"）送入神经网络？

        \item \textbf{特征标准化}：为什么要对数值特征进行 \texttt{StandardScaler} 标准化？如果房价数据中"面积"特征范围是 $[50, 500]$ 平方米，而"房间数"特征范围是 $[1, 10]$，不标准化会带来什么问题？

        \item \textbf{输出层设计}：模型的最后一层 \texttt{nn.Linear(128, 1)} 没有接激活函数，解释这是为什么。如果任务改成"预测房价是否超过500万（二分类）"，输出层应该怎么修改？

        \item \textbf{训练过程}：在一个完整的训练循环中，\texttt{optimizer.zero\_grad()}、\texttt{loss.backward()} 和 \texttt{optimizer.step()} 三个步骤必须按顺序执行，缺一不可。解释这三个调用各自的作用，以及为什么 \texttt{zero\_grad()} 必须在 \texttt{backward()} 之前调用。

        \item \textbf{评估指标}：MSE 损失虽然便于计算，但直觉上不好理解（单位是"元的平方"）。实际房价预测中通常报告哪种指标？写出其公式，并说明它相比 MSE 的优势。
    \end{enumerate}

    \vspace{0.2cm}
    {\color{gray}\small\textit{来源溯源：[文档第 7 章 7.6 应用案例：房价预测]}}
\end{problembox}

\begin{beginnerbox}[小白必读：训练一个神经网络，到底在做什么？]
    \begin{itemize}
        \item \textbf{一次训练迭代的四步节奏}：(1) 前向传播——数据流过网络，得到预测值；(2) 计算损失——预测值和真实值的差距；(3) 反向传播——自动计算所有参数的梯度；(4) 参数更新——优化器按梯度方向调整参数。
        \item \textbf{为什么要多次迭代？}一次迭代通常只用一小批数据（mini-batch），神经网络需要看成千上万次数据才能学会规律。每轮完整过一遍数据集叫一个 epoch，通常训练需要几十到几百个 epoch。
        \item \textbf{训练集、验证集、测试集的区别}：训练集用来更新参数；验证集用来调超参数（学习率、网络结构），判断是否过拟合；测试集只在最终评估时用一次，不参与任何调整。
    \end{itemize}
\end{beginnerbox}

\begin{metaphorbox}[生活比喻：训练神经网络就像备考]
    \begin{itemize}
        \item \textbf{训练集} = 你平时刷的题目，做完对答案，每次都学到东西（更新权重）；
        \item \textbf{验证集} = 模拟考，用来判断自己当前水平，决定要不要换一本参考书（调超参数）；
        \item \textbf{测试集} = 真正的考试，只能用一次，不能拿来练手；
        \item \textbf{过拟合} = 题目背得滚瓜烂熟，但换一道新题就抓瞎——模型在训练集上表现很好，在测试集上表现很差；
        \item \textbf{Dropout 和正则化} = 故意在复习时遮住部分笔记，逼自己真正理解，而不是死记硬背。
    \end{itemize}
\end{metaphorbox}

\begin{solutionbox}[参考答案与详细解析]
    \textbf{(1) 独热编码的原理及必要性：}

    独热编码将 $k$ 类的类别变量转换为 $k$ 维二进制向量，每次只有一个维度为1，其余为0。例如"北京=1，上海=2，广州=3"→"北京=[1,0,0]，上海=[0,1,0]，广州=[0,0,1]"。

    不能直接用整数编码的原因：神经网络会把整数值当作连续数值处理，隐含"北京离上海的距离=1，北京离广州的距离=2"这样的大小关系，而三个城市在特征语义上本无大小之分，这种人为引入的大小关系会误导模型学习错误的规律。

    \textbf{(2) 标准化的必要性：}

    梯度下降对各特征的尺度非常敏感。面积范围 $[50, 500]$ 和房间数范围 $[1, 10]$ 相差50倍，对应的梯度幅度差异巨大：面积特征上的梯度可能比房间数的梯度大50倍。这导致优化时只能使用极小的学习率（否则面积方向震荡），房间数方向更新极慢，整体收敛速度大幅下降，甚至难以收敛。标准化后两者都在均值0、方差1的尺度上，梯度幅度一致，可以使用更大的学习率，收敛更稳定。

    \textbf{(3) 输出层无激活函数的原因 \& 分类改造：}

    房价预测是回归任务，预测值应该是任意正实数（如200万、580万），若加 Sigmoid 输出只能在 $(0,1)$，若加 ReLU 又会引入非必要的截断。因此输出层直接保留线性输出，不加激活函数。

    改为二分类任务的修改：
    \begin{lstlisting}[language=Python]
# 修改方案一（推荐）：输出 logit，配合 BCEWithLogitsLoss
nn.Linear(128, 1)   # 输出层不加激活函数
loss_fn = nn.BCEWithLogitsLoss()  # 内置 Sigmoid

# 修改方案二：输出概率，配合 BCELoss
nn.Sequential(nn.Linear(128, 1), nn.Sigmoid())
loss_fn = nn.BCELoss()
    \end{lstlisting}

    \textbf{(4) 三个步骤的作用：}

    \begin{itemize}
        \item \textbf{\texttt{optimizer.zero\_grad()}}：清空所有参数上一次迭代积累的梯度。PyTorch 默认将梯度\textbf{累加}而不是覆盖，若不清零，本次梯度会叠加在历史梯度上，导致参数更新方向错误；
        \item \textbf{\texttt{loss.backward()}}：触发自动微分，从损失出发反向计算所有参数的梯度，存储在各参数的 \texttt{.grad} 属性中；
        \item \textbf{\texttt{optimizer.step()}}：读取各参数的 \texttt{.grad}，按优化算法（如 Adam）更新参数值。
    \end{itemize}

    \texttt{zero\_grad()} 必须在 \texttt{backward()} 之前，否则新的梯度会累加在旧梯度上，使得参数更新错误，训练无法正常进行。

    \textbf{(5) RMSE 评估指标：}

    实际中常用均方根误差（RMSE，Root Mean Squared Error）：
    \[
    \text{RMSE} = \sqrt{\frac{1}{N}\sum_{i=1}^{N}(\hat{y}_i - y_i)^2} = \sqrt{\text{MSE}}
    \]
    优势：单位与房价相同（元而非元的平方），直觉上易于理解——"模型预测误差平均约为 $X$ 元"。此外，有时还使用平均绝对百分比误差（MAPE）衡量相对误差：
    \[
    \text{MAPE} = \frac{1}{N}\sum_{i=1}^{N}\left|\frac{\hat{y}_i - y_i}{y_i}\right| \times 100\%
    \]

    \begin{keybox}[房价预测完整训练循环模板]
        \begin{lstlisting}[language=Python]
for epoch in range(num_epochs):
    model.train()                    # 切换训练模式
    for X_batch, y_batch in dataloader:
        optimizer.zero_grad()        # 第一步：清空梯度
        pred = model(X_batch)        # 第二步：前向传播
        loss = loss_fn(pred, y_batch)# 第三步：计算损失
        loss.backward()              # 第四步：反向传播
        optimizer.step()             # 第五步：更新参数

    model.eval()                     # 切换推理模式
    with torch.no_grad():            # 不计算梯度
        val_pred = model(X_val)
        val_loss = loss_fn(val_pred, y_val)
    print(f"Epoch {epoch}: val_loss={val_loss:.4f}")
        \end{lstlisting}
    \end{keybox}
\end{solutionbox}

\begin{appbox}[现实应用场景]
    \textbf{神经网络房价预测——实际中比想象中更复杂：}\\
    Zillow（美国最大房产平台）的 Zestimate 模型，就是用神经网络估算房价的商业落地。它面对的挑战包括：特征高达数百个（面积、朝向、学区评分、周边设施……），且各地区特征分布差异极大（北京的均价和县城均价不能用同一套权重）。实际中会用到更复杂的技巧，如地理位置嵌入（把经纬度变成向量）、特征交叉、对目标值取对数（$\log(\text{SalePrice})$）来压缩量纲。但本章案例的全流程骨架——特征工程 → 建模 → 选损失函数 → 训练 → 评估——与工业级方案完全一致。
\end{appbox}

% ===== 本章小结 =====
\section{本章小结：把零件装成一台可以跑的机器}

学完这 6 道题，如果还觉得"概念飘在空中"，试试把整章内容浓缩成这几句画面：

\begin{itemize}
    \item \textbf{激活函数}：ReLU 是隐藏层的标配；Sigmoid 输出概率（二分类）；Softmax 输出概率分布（多分类）；回归输出层什么都不加。
    \item \textbf{参数初始化}：绝对不能全初始化为同一个值；ReLU 配 Kaiming，Sigmoid/Tanh 配 Xavier。
    \item \textbf{损失函数}：分类用交叉熵，回归用 MSE/RMSE；\texttt{CrossEntropyLoss} 自带 Softmax，不要重复。
    \item \textbf{优化器}：大多数任务直接用 Adam；需要精调时，学习率衰减是让训练最终收敛的关键。
    \item \textbf{模型结构}：简单的用 \texttt{Sequential}，需要灵活结构的用 \texttt{nn.Module} 子类；\texttt{forward} 里追踪每层输出的形状是调试的第一步。
    \item \textbf{训练循环}：\texttt{zero\_grad → backward → step}，顺序不能乱；训练时 \texttt{model.train()}，推理时 \texttt{model.eval()}，别忘了切换。
\end{itemize}

\begin{metaphorbox}[给零基础读者的一句总比喻]
    \textbf{如果把训练一个神经网络比作培养一名厨师：}
    \begin{itemize}
        \item \textbf{激活函数} = 厨师的直觉（非线性判断力）：没有它，再多的训练也只会做直来直去的"线性菜"；
        \item \textbf{初始化} = 开业前的准备状态：从一个好的起点出发，比从乱糟糟的状态开始训练省很多时间；
        \item \textbf{损失函数} = 顾客的反馈表单：反馈维度选对了，厨师才知道往哪个方向改进；
        \item \textbf{优化器} = 厨师改进菜谱的方法论：是把每次差评都死记硬背（SGD），还是总结规律系统改进（Adam）？
        \item \textbf{最终效果}：一个调好参数、训练充分的神经网络，就像一位经过千百次反馈迭代、炉火纯青的大厨——它"记住"的不是具体的食谱，而是做好菜的底层规律。
    \end{itemize}
\end{metaphorbox}

\end{document}
